\documentclass[../main.tex]{subfiles}
\graphicspath{{\subfix{../images/}}}
\begin{document}
	Si può eseguire l'accesso ad una base di dati nei seguenti modi:
	\begin{itemize}
		\item Linguaggi testuali interattivi 
		\item Comandi inclusi in estensioni di linguaggi tradizionali 
		\item Comandi inclusi in linguaggi di sviluppo ad hoc 
		\item Interfacce grafiche amichevoli
	\end{itemize}
	\subsection{Linguaggi per basi di dati}
	\begin{itemize}
		\item Linguaggi di definizione dei dati\\
		Utilizzati per definire gli schemi e le autorizzazioni per l’accesso
		\item Linguaggi di manipolazione dei dati\\
		Utilizzati per l’interrogazione e l’aggiornamento dei contenuti della base di dati
	\end{itemize}
	Alcuni linguaggi specializzati, come ad esempio SQL, presentano le caratteristiche di entrambi i tipi di linguaggio.
	\subsubsection{SQL: un linguaggio interattivo}
	Come esempio immaginiamo di interfacciarci con la seguente base di dati:
	\begin{table}[ht]
		\begin{minipage}{.2\linewidth}
			\centering
			\begin{tabular}{ |c|c| }
				\hline
				\multicolumn{2}{|c|}{\textbf{Corsi}} \\
				\hline
				\textbf{Corso} & \textbf{Aula} \\
				\hline
				Sistemi & N3 \\
				Reti & N2\\
				\hline
			\end{tabular}
			%\caption{Caption}
		\end{minipage}
		\hfill
		\begin{minipage}{.2\linewidth}
			\centering
			\begin{tabular}{ |c|c| }
				\hline
				\multicolumn{2}{|c|}{\textbf{Aule}} \\
				\hline
				\textbf{Nome} & \textbf{Piano} \\
				\hline
				N3 & Terra \\
				N2 & Terra \\
				\hline
			\end{tabular}
			%\caption{Caption}
		\end{minipage}
		\hfill
		\begin{minipage}{.2\linewidth}
			\begin{tabular}{ |c|c|c| }
				\hline
				\textbf{Corso} & \textbf{Aula} & \textbf{Piano}\\
				\hline
				Sistemi & N3 & Terra \\
				Reti & N2 & Terra \\
				\hline
			\end{tabular}
		\end{minipage}
	\end{table}
	\begin{lstlisting}[language=SQL]
		SELECT Corso, Aula, Piano 
		FROM Aule, Corsi 
		WHERE Nome = Aula AND Piano = "Terra"
	\end{lstlisting}
	\paragraph{SQL embedded in Pascal code}
	Di seguito un esempio di SQL integrato in un linguaggi di alto livello:
	\begin{lstlisting}[language=Pascal]
		write('nome della citta''?'); 
		readln(citta); 
		EXEC SQL DECLARE P CURSOR FOR 
		SELECT NOME, REDDITO 
		FROM PERSONE 
		WHERE CITTA = :citta ;
		EXEC SQL OPEN P ; 
		EXEC SQL FETCH P INTO :nome, :reddito ; 
		while SQLCODE = 0 do begin 
		write('nome della persona:', nome, 'aumento?'); 
		readln(aumento); 
		EXEC SQL UPDATE PERSONE SET REDDITO = REDDITO + 
		:aumento 
		WHERE CURRENT OF P 
		EXEC SQL FETCH P INTO :nome, :reddito 
		end; 
		EXEC SQL CLOSE CURSOR P
	\end{lstlisting}
	\paragraph{SQP embedded in PHP code}
	Connessione a mySQL tramite API PHP:
	\begin{lstlisting}[language=PHP]
		<html> <body>
		<?php
		$user="root"; $password=""; $host="127.0.0.1"; $database="prova";
		mysql_connect($host,$user,$password);
		@mysql_select_db($database) or die("Unable to select database");
		$query="SELECT * from utenti";
		$res=mysql_query($query);
		mysql_close();
		
		$row = mysql_fetch_assoc($res);
		print '<table><tr>';
		foreach($row as $name => $value) {
			print "<th>$name</th>";
		}
		print '</tr>';
		
		while($row) {
			print '<tr>';
			foreach($row as $value) {
				print "<td>$value</td>";
			}
			print '</tr>';
			$row = mysql_fetch_assoc($res);
		}
		print '</table>';
		?> </body> </html>
	\end{lstlisting}
	\subsection{Interazione non testuale}
	Come interazioni non testuali si intendono generalmente i DBMS come ad esempio Microsoft Access:
	\begin{figure}[h]
		\centering
		\adjustbox{cfbox=white 5pt 0cm}{\includegraphics[scale=0.3]{example-image-golden}}
		\caption{Realizza l’approccio detto “Query by Example”}
		\label{fig:im-4}
	\end{figure}
	\subsubsection{Vantaggi dei DBMS}
	\begin{itemize}
		\item Disponibilità dei dati a tutta una comunità 
		\item Modello unificato e preciso della realtà di interesse 
		\item Controllo centralizzato dei dati 
		\item Condivisione 
		\item Indipendenza dei dati
	\end{itemize}
	\subsubsection{Vantaggi dei DBMS}
	\begin{itemize}
		\item Prodotti costosi, complessi, che richiedono investimenti in hardware, software, personale.
		\item Forniscono un numero elevato di servizi, in modo integrato e difficilmente scorporabile se le esigenze dell’utente sono inferiori alle caratteristiche offerte.
	\end{itemize}
\end{document}